% page 413 of the facsimile (image p-412 of the scan)
% block 165.1 x 242.0 mm ; left margin 36.9 mm
\pgc{15.240mm}{0.000mm}{
\l{\cel{56}405}
\l{}
\l{oleo.\cel{2}Substanco grasa, liquida, obtenita de vejetanto, animalo,}
\l{\cel{3}mineralo, uzata por nutro, medicino, lumizo, la arti,ed ula}
\l{\cel{3}ceremonii religiala. - DEFIS.}
\l{}
\l{\sou{olfaktar}.\cel{2}(\sou{trans.}) Perceptar la odoro di (ulo, ulu). - EFIS.}
\l{}
\l{\sou{oligarko}. Membro di guvernistaro oligarkiala. - DEFIRS.}
\l{}
\l{\sou{oligarkio}. Formo di guvernado, en qua la autoritato imperala}
\l{\cel{3}praktikesas da mikra quanto de civitani privilejiera. DEFIRS.}
\l{}
\l{\sou{olibano.}\cel{2}(\sou{farmacio)}. Rezino qua donas incenso. - EFIS.}
\l{}
\l{\sou{oligocena}. (\sou{geol.)}\cel{3}Dicesas pri grupo di tereno-strati terciara}
\l{\cel{3}(sur eoceno, sub-an mioceno). - DEFIS.}
\l{}
\l{\sou{olim}.\cel{2}En epoko qua forfluis de longe ante nun (quankam ne anti-}
\l{\cel{3}qua.)\cel{2}L.}
\l{}
\l{\sou{olimpiado}. (\sou{antique, en Grekia).}\cel{1}Periodo de quar yari, de celebro}
\l{\cel{3}di la Olimpio-ludi til la celebro +nexta. - DEFIRS.}
\l{}
\l{\sou{olivo}. (\sou{bot.)}\cel{3}Frukto kernoza, elipsoida, verdatra, quan produk-}
\l{\cel{3}tas arboro ye foliaro nevelkebla, L.\cel{1}\sou{olea europaea,}\cel{1}e de qua}
\l{\cel{3}onu extraktas oleo manjebla. - DEFIRS.}
\l{}
\l{\sou{ombro}. I. Parto di surfaco lumizita, ube la lumo-radii intercept-}
\l{\cel{3}esas da korpo opaka. - II. La parto nelumizita qua, en surfaco}
\l{\cel{3}lumizata, reproduktas la formo di korpo opaka qua interceptas}
\l{\cel{3}la lumo-radii. EFIS.}
\l{}
\l{\sou{omento}. (\sou{kirurgio)}. Granda rifalduro di la peritonio, qua fluk-}
\l{\cel{3}tuas avan la intestino tenua. - DEFIRS.}
\l{}
\l{\sou{omisar}. (\sou{trans.)}\cel{3}Ne inkluzar, sive vole, sive pro oblivio,}
\l{\cel{3}(ulo quon onu devas dicar od agar). - EFIS.}
\l{}
\l{\sou{omleto}.\cel{2}(\sou{koquarto}).\cel{2}Disho ek ovi plaudita e koquita en padelo,}
\l{\cel{3}kun butro e kondimenti o spici. - DEFR.}
\l{}
\l{\sou{omna}. I. (\sou{avan substantivo plurala}).\cel{2}Ici ed iti, sen ecepto. -}
\l{\cel{3}II. (\sou{avan substantivo singulara)}.\cel{3}Ica samagrade kam ita, sen}
\l{\cel{3}dicerno nek distingo. - EIL.}
\l{}
\l{\sou{omnibuso.}\cel{2}Veturo publika, qua cirkulas en la quarteri diversa}
\l{\cel{3}di urbo, e haltas dum sua parkuro por lasar enirar (embarkar)}
\l{\cel{3}od ekirar (desembarkar) la voyajanti quin lu transportas bloke}
\l{\cel{3}po preco mikra. - DEFIRS.}
\l{}
\l{\sou{omnivora}. Qua su nutras ye nutrivi irgaspeca - devenanta de}
\l{\cel{3}animali o de vejetanti. - DEFIS.}
\l{}
\l{\sou{on(u)}. Pronomo personala nedefinita, uzata por indikar un o}
\l{\cel{3}plura personi ye maniero tre generala. - EF.}
}
